\scanpage{60}
aux espaces (plutôt, germes d'espaces ou topos) \deleted{$\unsure{V_{ij}}$} plus généraux que les $V_{ij}^*$ (et les $V_{ij}$), ce sont les suivants : on considère
\[
i\leq j\leq k\qquad\text{d'où}\qquad X_i\subset X_j\subset X_k
\]
on prend
\[
(10)\qquad V_{i,k}^j\overset{\mathrm{def}}=\overline V_{X_i,X_k}\setminus(R_{\deleted{\unsure{j}}\inserted{i},k}\cup\unsure{S_{j,k}|\Sigma_{j,k}})\ \deleted{\unsure{m}}
\]
où $R_{i,k}$ est défini dans (3) (3bis), et où
\[
(11)\qquad S_{j,k}\overset{\mathrm{def}}=\bigcup_{\substack{l\in I\\\text{avec}\\j\leq l<k}}X_l
\]
\margin{(\uline{NB} $S_{j,k}=\varnothing$ si $j=k$\\
$S_{j,k}=X_j$ si $k$ succ. de $j$)}

On a, en comparant avec (3bis), une \deleted{\unsure{partition}} \inserted{\unsure{partition}} de l'ens. $I^{(k)}\overset{\mathrm{def}}=\{l\in I\mid\unsure{l<k|l\leq k}\}$ en deux parties,
\[
\unsure{I^{*(k)'}|I^{(k)'}}=\{l\in I^{(k)}\mid X_l\not\supset X_j\}\quad\text{et}\quad
\unsure{I^{(k)''}|I^{(k)'}}=\{l\in I^{(k)}\mid X_l\supset X_j\},
\]
et on a
\[
(12)\qquad\left\{\begin{aligned}
R_{j,k}&=\bigcup_{l\in I^{(k)'}}X_l\\
S_{j,k}&=\bigcup_{l\in I^{(k)''}}X_l,
\end{aligned}\right.
\qquad\begin{gathered}
\text{Donc }R_{j,k}\cup S_{j,k}=\dot X_k\\
(\text{aussi }R_{j,k}\cap S_{j,k}=\dot X_{\unsure{i|j}})
\end{gathered}
\]
\deleted{$\unsure{V_{ij}}$} On a donc
\[
(13)\qquad V_{i,k}^i=V_{i,k}^*,\qquad V_{i,k}^k=V_{i,k},
\]
\margin{d'où $V_{i,i}^i=V_{i,i}=V_{i,i}^*=X_i^*$.}

$V_{i,k}^j$ a encore tendance à être \deleted{\unsure{un}} homéomorphe (ou du moins, \unsure{homotope}, ou homéomorphe en un sens affaibli \unsure{de germes|de genre}, qu'il faudrait dégager\ldots) à un fibré loc. trivial de base $X_i^*$, mais comme c'est déjà dans le cas $V_{i,k}^k=V_{i,k}$) à fibre pouvant être \uline{singulière}. Dans le cas de $V_{i,k}^k$,
\mspage{59}
